\documentclass{article}
\usepackage[T1]{fontenc}
\usepackage{lmodern}
\usepackage{graphicx}
\usepackage{amsmath,amssymb}
\usepackage[a4paper,margin=2cm]{geometry}
\usepackage[absolute,overlay]{textpos}
\setlength{\TPHorizModule}{1mm}
\setlength{\TPVertModule}{1mm}
\usepackage[indonesian]{babel}
\usepackage{hyperref}
\setlength{\emergencystretch}{2em}
% unit: Halaman 65

\begin{document}

% === Halaman 65 (python/native) ===

Jaringan Feistel

• Beberapa cipher blok seperti DES, GOST, Feal, Lucifer, RC5, RC6, dan 

lain-lain menggunakan struktur enciphering pada setiap putaran yang 

dinamakan jaringan Feistel (Feistel Network).

• Dalam hal ini, blok plainteks dibagi menjadi dua bagian, dan pada 

masing-masing bagian dilakukan proses transformasi menjadi upa-

bagian pada putaran selanjutnya. 

• Struktur ini bersifat reversible, yaitu untuk melakukan dekripsi maka 

jaringan Feistel dieksekusi dari “bawah” ke “atas”. 

65

\end{document}
